\evidencelead{E01 · 用户纠正 / 原始对话}{不只缩减Goal数量，更要改变执行机制}{录屏约00:48—01:18：先保留我的原话，再看GPT如何将其转成研究合同。}
\evidenceshot{D01_user_goal}{我的第一项纠正}{我要求最多两个、最好一个Goal，并指出“做完就停”和缺少可视化研发的问题。}{16}{119}{24}{C2}
\sourceband{RAW VIDEO 00:48.500 · E01 / USER · 原句未经改写，截图末尾可见原生“显示更多”控件}
\noindent 我不希望只有“先把TODO跑通，然后给出一个完成状态”。这条消息还要求从多个简单实验入手，避免一开始混合策略时无法定位失败原因。GPT在随后回复中重新区分了完成、研究与优化。

\evidencelead{E01 · 同一互动续页}{我批准的不是无界自动搜索}{同一条用户消息的后半段；它说明长任务自主权成立的前提。}
\evidenceshot{D01_user_longrun}{授权需要边界}{我提出用真实结果持续扩展控制变量实验，允许有依据的组合；只有明确授权长任务时，Codex才可自主推进。}{44}{120}{70}{C2}
\sourceband{RAW VIDEO 00:51.500 · E01 / USER CONTINUATION · 与上页同一消息，未重排成新用户轮次}
\evidenceshot{D01_gpt_onegoal}{GPT的实质修订}{GPT将方案收束为一个父Goal，并明确研究合同、运行中优化与最终验收分离。}{27}{119}{59}{C1}
\sourceband{RAW VIDEO 00:55.500 · E01 / GPT AFTER · “一个完整的长Goal”}\par
